% Source: https://artofproblemsolving.com/community/c3h2934037_2015_emcc_guts_round_p13_p24__exeter_math_club_competition
% Stable identity: aops-topic-2934037-q24
Let $v,e$ be the vertex and edge counts. Then $3v=2e$, so $v$ is even. A connected component has at least 4 vertices. With 4 vertices the graph would be $K_4$, impossible for unit distances: a unit equilateral triangle has no point at unit distance from all three vertices. With 6 vertices the complement is a 2-regular simple graph, hence either two triangles or a 6-cycle. The original graph is respectively $K_{3,3}$ (nonplanar) or the triangular prism. The prism cannot have a planar unit-edge drawing. If its outer face is a triangle, that unit equilateral triangle cannot contain the other disjoint unit triangle of equal area. Otherwise the outer face is a quadrilateral, hence a rhombus. Its three quadrilateral faces have corresponding opposite edges parallel, so the two triangular faces are translates by a vector $w$. For the three oriented sides $u_1,u_2,u_3$ of one triangle, $u_1+u_2+u_3=0$. The three rhombus areas are $|u_i\times w|$; each is at most the sum of the other two. But the outer rhombus must contain the two other rhombi and two positive-area triangles, a contradiction. Thus every possible component has at least 8 vertices, so $e\ge12$. For attainment, form a diamond from the two equilateral triangles on the edge joining $(-\dfrac{1}{2},0)$ to $(\dfrac{1}{2},0)$; its other vertices are $(0,\pm\dfrac{\sqrt3}{2})$. Translate this five-edge diamond by $(\dfrac{\sqrt3}{2},\dfrac{1}{2})$, and join each pair of corresponding degree-two vertices. The two added segments have length 1, and the displayed coordinates give disjoint diamonds and noncrossing joining edges. All eight vertices have degree 3. Hence $\boxed{12}$.
\begin{asy}
usepackage("amsmath");
size(7cm);
pair A=(-.5,0),B=(.5,0),C=(0,sqrt(3)/2),D=(0,-sqrt(3)/2),v=(sqrt(3)/2,.5);draw(A--C--B--D--cycle);draw(A--B);draw((A+v)--(C+v)--(B+v)--(D+v)--cycle);draw((A+v)--(B+v));draw(C--(C+v)^^D--(D+v),blue);pair[] p={A,B,C,D,A+v,B+v,C+v,D+v};for(int i=0;i<8;++i){dot(p[i]);label(string(i+1),p[i],i%2==0?NW:SE);}label("12 unit segments",(.4,-1.3));
\end{asy}

\begin{asy}
usepackage("amsmath");
size(7cm);
pair A=(0,0),B=(4,0),C=(2,3),D=(1,1),Ee=(3,1),F=(2,2);draw(A--B--C--cycle);draw(D--Ee--F--cycle);draw(A--D^^B--Ee^^C--F);label("outer triangle",(2,-.5));label("inner triangle",(2,1.3));label("schematic prism case; lengths not asserted",(2,-1.1));
\end{asy}

\begin{asy}
usepackage("amsmath");
size(7cm);
pair A=(0,0),B=(4,0),Ee=(4,4),D=(0,4),C=(2,1.2),F=(2,2.8);draw(A--B--Ee--D--cycle);draw(A--C--B^^D--F--Ee^^C--F);label("triangle",(2,.4));label("triangle",(2,3.6));label("$R_1$",(.7,2));label("$R_2$",(3.3,2));label("outer quadrilateral",(2,-.5));label("schematic prism case; lengths not asserted",(2,-1.1));
\end{asy}
